\section{Convexity on the feasible domain: a cubic--quartic boundary}
\label{sec:domain-boundary}

Global convexity is a substantive assumption in the preceding point
theorem. Convexity only on a polytope still suffices for cubics, with
effective constants and a fixed selector. For quartics, even a constant
accuracy point query can decide difficult arithmetic comparisons. These
statements concern the explicitly encoded polynomial and the output
contracts of Section~\ref{sec:contracts}; they do not assert that the
underlying arithmetic problems are outside polynomial time.

\subsection{Effective point approximation for cubics on polytopes}

\begin{theorem}[Domain-convex cubic point oracle]
\label{thm:cubic-domain}
Let $P=\{x\in\R^n:Cx\le b\}$ be a nonempty bounded rational
polytope. Let $f\in\Q[x]$ be an explicitly supplied polynomial of
degree at most three, promised convex on $P$. If $I$ is the total
binary input length, a deterministic algorithm computes a rational
$\Gamma_P\ge1$ with $\poly(I)$ bits such that, for
$S=\argmin_P f$,
\begin{equation}
 \dist(x,S)\le\Gamma_P\bigl(f(x)-f_*\bigr)^{1/4}
 \quad\text{when }x\in P,\quad0\le f(x)-f_*\le1.
 \label{eq:cubic-error}
\end{equation}
For every $q\ge0$, it returns a feasible rational point within
$2^{-q}$ of $S$ in time $\poly(I+q)$. It can instead approximate the
unique minimum-Euclidean-norm point of $S$ to the same accuracy and in
the same time bound. The norm is taken in the original coordinates.
No uniqueness, strong convexity, or optimizer-interiority assumption
is required. Convexity recognition is outside the theorem.
\end{theorem}

The proof is included at the level needed to identify the effective
constants and the role of degree three. The complete coefficient
estimates and weak-optimization interface are in the
\href{../../research-20261002/new-direction/convex-cubic-polytope-point-oracle.md}{cubic polytope companion proof}.

\paragraph{Affine hull and a rational interior ball.}
For each row of $Cx\le b$, rational linear programming determines
whether its maximum slack on $P$ is zero. These universally tight
rows define $\operatorname{aff}(P)$: averaging a strict-slack witness
for every other row gives a point strict in all of those rows at once.
Rational elimination therefore gives
\[
 x=x_0+Vu,\qquad P'=\{u:Au\le d\},\qquad
 \phi(u)=f(x_0+Vu),
\]
with $P'$ full dimensional and bounded. Choose free original
coordinates as $u$, so that $V$ contains an identity submatrix.
All these data have polynomial binary length. Explicit substitution
preserves polynomial size because the degree is at most three.
The zero-dimensional case is immediate.

For each nonzero row $a_i^{\mathsf T}$ of $A$, impose
$a_i^{\mathsf T}c+\rho\|a_i\|_1\le d_i$, and maximize
$\rho$ subject to $0\le\rho\le1$. Full dimension gives a positive
rational optimum. Thus $B(c,\rho)\subseteq P'$. Coordinate linear
programs also give a rational $R\ge1$ with
$\|u-c\|\le R$ on $P'$. The logarithms of $R$ and $1/\rho$
are polynomially bounded. This reduction must precede the Hessian
argument: convexity on a lower-dimensional set need not imply
positive semidefiniteness of the ambient Hessian.

\paragraph{An affine Hessian identifies the relevant directions.}
Write $H(u)=\nabla^2\phi(u)$ and $H_c=H(c)$. The matrix $H(u)$
is affine in $u$ and positive semidefinite on $P'$. For each
$u\in P'$, the reflected point
$u'=c-(\rho/R)(u-c)$ lies in the interior ball. Taking the affine
combination that equals $c$ gives
\begin{equation}
 0\preceq H(u)\preceq\beta H_c,\qquad
 \beta=1+R/\rho.
 \label{eq:cubic-domination}
\end{equation}
Consequently
$K=\ker H_c=\bigcap_{u\in P'}\ker H(u)$ is a known rational
subspace. Boundary Hessians may have larger kernels. Choose a rational
$M\ge1$ bounding $\|H(u)\|$ on $P'$; for example,
$M=\max\{1,\beta\max_i\sum_j|(H_c)_{ij}|\}$ works.
If $H_c=0$, the objective is affine and linear programming handles
optimization; the same linear error-bound argument below supplies
\eqref{eq:cubic-error} in this branch.

Let $v$ be any optimizer of $\phi$, $u\in P'$, $d=u-v$, and
$E=\phi(u)-\phi(v)$. Put
$a=d^{\mathsf T}H(v)d$ and $b=d^{\mathsf T}H(u)d$.
Exact cubic Taylor expansion and first-order optimality imply
\[
 E=\nabla\phi(v)^{\mathsf T}d+(2a+b)/6\ge(a+b)/6.
\]
Full symmetry of the constant third derivative gives the additional
identity
\[
 d^{\mathsf T}H_cd
 =a+(c-v)^{\mathsf T}(H(u)-H(v))d.
\]
For a positive semidefinite matrix $Q$,
$\|Qd\|\le\sqrt{\|Q\|d^{\mathsf T}Qd}$.
Combining this inequality with the diameter bound $2R$ gives the
conservative estimate
\begin{equation}
 E\ge\frac{(d^{\mathsf T}H_cd)^2}{384MR^2}
 \ge\frac{\lambda^2\|\Pi d\|^4}{384MR^2},
 \label{eq:cubic-transverse}
\end{equation}
where $\Pi$ projects onto $K^\perp$ and $\lambda$ is the smallest
positive eigenvalue of $H_c$. The unknown optimizer and this projector
are used only in the analysis.

\paragraph{The optimizer set is an affine slice.}
Put $g_c=\nabla\phi(c)$. Integrating the Hessian shows that
$\nabla(\phi-g_c^{\mathsf T}u)$ belongs to $K^\perp$ throughout
$P'$. Equation~\eqref{eq:cubic-transverse} then proves
\begin{equation}
 \argmin_{P'}\phi
 =\{w\in P':H_c(w-v)=0,\ g_c^{\mathsf T}(w-v)=0\}.
 \label{eq:cubic-slice}
\end{equation}
The gradient row is needed because $\phi$ can have a nonzero affine
slope along the common Hessian kernel. It satisfies
\[
 |g_c^{\mathsf T}d|\le E+MR\|\Pi d\|.
\]
The slice right-hand side may be irrational; its coefficient matrix
is known and rational.

Here is the quantitative linear error bound used for that slice.
If $Q=\{u:Bu\le h\}$ and
$Z=\{u\in Q:Tu=t\}\ne\varnothing$, with integer matrices
$B,T$ having $s$ columns and entries bounded in absolute value by
$C_*\ge1$, then
\begin{equation}
 \dist(u,Z)\le(sC_*)^{s-1}\|Tu-t\|\qquad(u\in Q).
 \label{eq:integer-hoffman-domain}
\end{equation}
Neither right-hand side has to be rational. To see the stated height
bound, project $u$ onto $Z$ and express the displacement using
independent equality normals and a nonnegative combination of active
inequality normals. Eliminate dependence modulo the equality span.
The resulting integer row matrix has at most $s$ independent rows;
its Gram determinant is at least one and its singular values are at
most $sC_*$. Its smallest singular value is therefore at least
$(sC_*)^{-(s-1)}$. Taking the inner product with the displacement
discards the nonpositive contributions of active outward inequality
normals, since $u\in Q$, and proves
\eqref{eq:integer-hoffman-domain}.

Clear denominators of $H_c,g_c,A$ by a positive integer $\Delta$.
If $H_c$ has positive rank, a nonzero principal minor gives the
rational eigenvalue lower bound
$\lambda\ge\Delta^{-s}M^{-(s-1)}$.
Equations~\eqref{eq:cubic-transverse}--\eqref{eq:integer-hoffman-domain}
now give a computable constant of polynomial binary length multiplying
$E^{1/4}$ for $E\le1$. Multiplication by a rational upper bound on
$\|V\|$ proves \eqref{eq:cubic-error} in original coordinates.
When $H_c=0$, the gradient row alone gives an error bound proportional
to $E$, hence also to $E^{1/4}$ on this range.

\paragraph{Feasible output and a consistent selector.}
Classical weak convex optimization applies to the capped epigraph with
the known inner and outer radii \citep{GroetschelLovaszSchrijver1988}. Approximate
feasibility must be repaired. If a rational $u$ is within $\delta$
of $P'$, the rational point
\[
 \widehat u=\frac{\rho u+\delta c}{\rho+\delta}
\]
lies in $P'$ exactly. Relative to a nearby feasible point the repair
costs at most $\delta(1+R/\rho)$. A coefficient bound on
$\|\nabla\phi\|$ controls its objective cost. Together with the
usual erosion bound for the capped epigraph, this yields certified
value intervals of width $\eta$ in time
$\poly(I+\log(1/\eta))$. Choosing
$\eta=(2^{-q}/\Gamma_P)^4$ proves the set-distance statement.

For completeness, let $p$ be the minimum-original-norm optimizer and
let $R_x\ge1$ bound $\|x\|$ on $P$. Set
\[
 \varepsilon=2^{-q},\qquad
 \tau=\frac{\varepsilon^6}{1024R_x^4\Gamma_P^4},\qquad
 \eta=\tau\varepsilon^2/4.
\]
Minimize $f(x)+\tau\|x\|^2$ to feasible value gap $\eta$.
If $x_\tau$ is its exact minimizer and
$e=\dist(x_\tau,S)$, projection onto the closed convex set $S$
and regularized optimality give
\[
 e^4/\Gamma_P^4\le f(x_\tau)-f_*
 \le2R_x\tau e,\qquad
 \|x_\tau-p\|^2\le2R_xe.
\]
The gap is at most $\tau R_x^2\le1$, so the error bound is
applicable. These inequalities give selection bias at most
$\varepsilon/2$. Strong convexity of the regularized objective gives
another $\varepsilon/2$ for its value approximation. In reduced
coordinates the penalty is $\tau\|x_0+Vu\|^2$; replacing it with
$\tau\|u\|^2$ could select a different optimizer.

The exponent $1/4$ is sufficient for polynomial dependence on
precision bits. The example $f(t)=t^3$ on $[0,1]$ excludes a uniform
error exponent larger than $1/3$. The present argument does not prove
that $1/3$ is attainable with effective constants in all dimensions.
That sharper exponent is not needed for Theorem~\ref{thm:cubic-domain}.

\subsection{Quartics can encode an arithmetic sign in one endpoint}

\begin{theorem}[Constant-accuracy point reductions]
\label{thm:quartic-point-lower}
There are polynomial-time reductions with one point-oracle query from
each of the following decision problems to finding a point within
Euclidean distance $1/4$ of the optimizer of an explicitly encoded
rational quartic convex on a rational box:
\begin{enumerate}[label=(\roman*)]
\item Square Root Sum, with positive integer radicands and an integer
threshold. The output instance has interaction-treewidth at most two,
bounded coefficient magnitudes, and bounded upper coordinate curvature.
\item $\PosSLP$. The output instance has polynomial encoding length
and can have bounded coefficient magnitudes; no bounded-treewidth
claim is made.
\end{enumerate}
Both instances can be written on unit boxes and have unique
optimizers. In both, a designated optimizer coordinate equals zero or
one according to the source comparison. Neither reduction supplies a
uniform effective lower growth bound for the final objective.
\end{theorem}

The complete constructions are the
\href{../../research-20261002/new-direction/convex-point-radical-comparison.md}{radical point reduction}
and the
\href{../../research-20261002/new-direction/posslp-convex-point-extraction.md}{arithmetic-circuit point reduction}.
The following common step explains why the accuracy remains constant
even when the arithmetic sign has extremely small magnitude.

Suppose two independent strongly convex base problems contain
coordinates $a_+,a_-\in[0,1]$, with exactly one positive at their
respective unique minima. The positive coordinate records which side
of a strict comparison holds. With a fresh $y\in[0,1]$, add
\begin{equation}
 \eta\bigl[a_+^2(y-1)^2+a_-^2y^2\bigr],
 \qquad\eta=1/192.
 \label{eq:endpoint-amplifier}
\end{equation}
For a direction $(p,q)$ in $(a,y)$, direct differentiation gives
\[
 D^2\bigl[a^2(y-c)^2\bigr][(p,q)]^2
 =2\bigl(aq+2(y-c)p\bigr)^2-6(y-c)^2p^2.
\]
The possible negative contributions are absorbed by the base
curvature, which is at least $1/16$ in the radical construction and
at least $1/2$ in the unscaled circuit construction. This proves
joint convexity on the box, including where an amplitude vanishes.

Every full objective is bounded below by the sum of its two base
minima. If $a_+^*>0$, choosing both base minimizers and $y=1$
attains this bound; if $a_-^*>0$, choose $y=0$. Every minimizer
must therefore attain both unique base minima and make
\eqref{eq:endpoint-amplifier} zero. The full optimizer is unique
and its endpoint is forced. Thresholding an approximation to $y$ at
$1/2$ decides the source comparison. Feasibility of that approximation
is not needed for the implication.

In the radical reduction, normalized cubic leaf objectives have
minimizers $\sqrt{a_i}/R_i\in[1,2]$, and a binary averaging tree
propagates their weighted sum. Opposite linear couplings to the root
produce the two amplitude tests. Equality is removed by polynomial-time
integer square-root tests: a positive sum of square roots of integers
can be an integer only when all its summands are integers. One proof
takes traces in the common multiquadratic extension; every nonsquare
radical has trace zero, and positivity excludes their remaining sum
being zero. The algorithm does not construct that extension or factor
the radicands. The tree factors need bags of size at most three, and
the two amplitude--endpoint edges join the two trees without increasing
that bound. After rescaling to the unit box, remove the aggregate constant
term; this preserves the minimizer and the bounded-coefficient promise.

In the circuit reduction, bounded rational-pair gate values encode the
integer circuit without expanding its potentially enormous output.
Triangular degree-two residual equations have a unique common zero.
Geometrically decreasing weights on their squares dominate the negative
Hessian terms and produce a quartic strongly convex on its box.
The final signal has the sign of $2A-1$, where $A$ is the integer
source output, so it never vanishes. Opposite amplitude tests and
\eqref{eq:endpoint-amplifier} finish the reduction. The companion proof
checks arbitrary fanout and repeated operands. Scaling the entire
objective makes its collected coefficient magnitudes bounded while
preserving its minimizer; this also scales its curvature.

Before optional whole-objective scaling, the endpoint curvature is
$2\eta((a_+^*)^2+(a_-^*)^2)$ and can be extremely small. Thus a
value-gap guarantee does not give the requested point guarantee.
Global convexity is also absent: box convexity is exactly what the
Hessian estimates establish. These qualifications explain compatibility
with the positive global-convex and cubic-domain results.

An ordinary polynomial-time point algorithm for either target class
would put its source problem in P. An always-correct algorithm with
expected polynomial work would give the corresponding Las Vegas
algorithm. These are conditional implications, not NP-hardness claims
or proved class separations. The distinction between approximate
values and strong approximation of solution coordinates has established
precedents in algebraic fixed-point complexity \citep{EY2010}; the
specialized convex-polynomial realizations require the separate proofs
cited above. Their independent
\href{../reviews/point-core/hardness-audit.md}{point-reduction audit}
and the
\href{../reviews/point-core/review.md}{positive-core audit}
record the checked assumptions. Neither audit establishes publication
priority.
